% ============================================================================
% SOLUTION-LAYER ARGUMENTATION MODULE (standalone)
% Sustainable University IT — Technical and Operational Objections and Responses
% Defends: ../solution-EN.tex
%          "A Technical Reference for Architecture and Migration"
%
% Scope: TECHNICAL / OPERATIONAL objections only —
%   architecture complexity, open-source coverage gaps (iOS/Android MDM, ...),
%   sizing realism, brownfield migration and coexistence, vendor lock-in vs
%   replaceability, conformance-test rigour, observability/audit load,
%   reversibility.
%
% Method: questionbox / answerbox / clarificationbox environments,
%   concession-then-rebuttal pattern, and the FIVE-LAYER evidence
%   methodology (legal > standards > sector > empirical > institutional)
%   with citations from the shared bibliography. Fully GENERIC.
% ============================================================================
\documentclass[11pt,a4paper,twoside]{book}

% === Encoding and language ===
\usepackage[utf8]{inputenc}
\usepackage[T1]{fontenc}
\usepackage[english]{babel}

% === Page layout ===
\usepackage[top=2.5cm,bottom=1.75cm,left=2.5cm,right=2.5cm]{geometry}
\usepackage{setspace}
\setstretch{1.0}

% === Typography and fonts ===
\usepackage{lmodern}
\usepackage{microtype}

% === Graphics and colours ===
\usepackage{graphicx}
\usepackage{suit-logo}
\usepackage[dvipsnames,table]{xcolor}
\definecolor{Slate}{RGB}{70,70,90}
\definecolor{Crimson}{RGB}{220,20,60}

% === Tables ===
\usepackage{booktabs}
\usepackage{tabularx}
\usepackage{longtable}
\usepackage{multirow}
\usepackage{array}

% === Lists ===
\usepackage{enumitem}

% === Symbols (checkmarks, crosses) ===
\usepackage{pifont}

% === Links ===
\usepackage[colorlinks=true,linkcolor=NavyBlue,citecolor=OliveGreen,
            urlcolor=BrickRed,plainpages=false,pdfpagelabels]{hyperref}
\usepackage{bookmark}
\bookmarksetup{numbered}
\usepackage{xurl}

% === Fix bookmark level: \section is used directly under \part (no \chapter) ===
\makeatletter
\renewcommand{\toclevel@section}{1}
\renewcommand{\toclevel@subsection}{2}
\renewcommand{\toclevel@subsubsection}{3}
\makeatother

% === Bibliography ===
\usepackage[backend=biber,
            style=numeric,
            sorting=none,
            sortcites=true,
            maxbibnames=99]{biblatex}
\addbibresource{../../shared/suit-consolidated.bib}

% === Citation display mode ===
\usepackage{bib-modes}
\setbibmode{linked}

% === natbib compatibility ===
\let\citep\parencite
\let\citet\textcite

% === Headers ===
\usepackage{fancyhdr}
\setlength{\headheight}{15pt}
\pagestyle{fancy}
\fancyhf{}
\fancyhead[LE,RO]{\thepage}
\fancyhead[RE]{\textit{Sustainable University IT --- Technical Objections \& Responses}}
\fancyhead[LO]{\textit{\leftmark}}
\fancyfoot[C]{\suitfootlogo}
\renewcommand{\headrulewidth}{0.4pt}
\renewcommand{\footrulewidth}{0.4pt}

% === Miscellaneous ===
\usepackage{csquotes}
\usepackage{float}
\usepackage{caption}
\usepackage{tcolorbox}
\tcbuselibrary{skins,breakable}

% === Glossary with tooltips ===
\usepackage{gls-modes}
% ============================================================================
% Glossary definitions for gls-modes.sty
% \glsdefine{key}{SHORT}{Long Form}{Tooltip definition}
% ============================================================================
%
% SUIT NOTE: this is the institution-neutral SUIT glossary. All entries are
% role-based / generic. Adopting institutions supply their own local values
% via the Localization Guide (shared/localization-guide-EN.tex) rather than
% via hard-coded glossary entries.
% ============================================================================

% --- Technical Infrastructure ---
\glsdefine{acl}{ACL}{Access Control List}{Access Control List -- a set of rules on a network device that permits or denies traffic based on source/destination addresses and ports.}
\glsdefine{api}{API}{Application Programming Interface}{Application Programming Interface -- a standardised interface allowing software systems to communicate with each other.}
\glsdefine{byod}{BYOD}{Bring Your Own Device}{Bring Your Own Device -- the practice of employees or students using personal devices to access institutional resources.}
\glsdefine{dhcp}{DHCP}{Dynamic Host Configuration Protocol}{Dynamic Host Configuration Protocol -- automatically assigns IP addresses to devices on a network.}
\glsdefine{dmz}{DMZ}{Demilitarised Zone}{Demilitarised Zone -- a network segment between an internal network and the Internet, hosting externally accessible services.}
\glsdefine{dns}{DNS}{Domain Name System}{Domain Name System -- translates human-readable domain names into IP addresses.}
\glsdefine{dpi}{DPI}{Deep Packet Inspection}{Deep Packet Inspection -- a method of examining the content of network packets as they pass through a checkpoint.}
\glsdefine{dtn}{DTN}{Data Transfer Node}{Data Transfer Node -- a server optimised for high-throughput scientific data transfers in a Science DMZ.}
\glsdefine{edr}{EDR}{Endpoint Detection and Response}{Endpoint Detection and Response -- security software monitoring endpoints for threats.}
\glsdefine{hpc}{HPC}{High-Performance Computing}{High-Performance Computing -- computing systems designed for large-scale scientific simulations and data processing.}
\glsdefine{iam}{IAM}{Identity and Access Management}{Identity and Access Management -- systems and processes for managing digital identities and controlling access to resources.}
\glsdefine{nac}{NAC}{Network Access Control}{Network Access Control -- system enforcing authentication, device posture assessment, and dynamic VLAN assignment at the network edge (e.g. PacketFence, Cisco ISE).}
\glsdefine{ids}{IDS}{Intrusion Detection System}{Intrusion Detection System -- monitors network traffic for suspicious activity (detection only, no inline blocking).}
\glsdefine{ips}{IPS}{Intrusion Prevention System}{Intrusion Prevention System -- monitors and actively blocks detected threats inline.}
\glsdefine{lms}{LMS}{Learning Management System}{Learning Management System -- platform for delivering and managing educational content (e.g. Moodle).}
\glsdefine{mfa}{MFA}{Multi-Factor Authentication}{Multi-Factor Authentication -- requiring two or more verification factors to access a resource.}
\glsdefine{pki}{PKI}{Public Key Infrastructure}{Public Key Infrastructure -- framework for managing digital certificates and encryption keys.}
\glsdefine{rbac}{RBAC}{Role-Based Access Control}{Role-Based Access Control -- access permissions assigned to roles rather than to individual users.}
\glsdefine{sciencedmz}{Science DMZ}{Science Demilitarised Zone}{Science DMZ -- a network architecture pattern developed by ESnet for high-performance scientific data transfers, bypassing enterprise firewalls in favour of ACLs and dedicated IDS.}
\glsdefine{sdn}{SDN}{Software-Defined Networking}{Software-Defined Networking -- centralised, programmable control of network behaviour.}
\glsdefine{siem}{SIEM}{Security Information and Event Management}{Security Information and Event Management -- centralised platform for collecting and analysing security logs and alerts.}
\glsdefine{soc}{SOC}{Security Operations Centre}{Security Operations Centre -- team and facility responsible for monitoring and responding to security incidents.}
\glsdefine{sso}{SSO}{Single Sign-On}{Single Sign-On -- authentication allowing access to multiple systems with one set of credentials.}
\glsdefine{tac}{TAC}{Trusted Activity Context}{Trusted Activity Context -- a coherent, validated environment assembling data, tools, services, and access rights for a specific university activity (concept introduced in this report).}
\glsdefine{vdi}{VDI}{Virtual Desktop Infrastructure}{Virtual Desktop Infrastructure -- technology delivering desktop environments from a centralised server to any device.}
\glsdefine{vlan}{VLAN}{Virtual Local Area Network}{Virtual Local Area Network -- a logical subdivision of a physical network, isolating traffic between groups of devices.}
\glsdefine{vpn}{VPN}{Virtual Private Network}{Virtual Private Network -- encrypted tunnel for secure remote connectivity.}

% --- Governance Frameworks and Standards ---
\glsdefine{abac}{ABAC}{Attribute-Based Access Control}{Attribute-Based Access Control -- access decisions based on attributes of the user, resource, and environment (NIST SP 800-162).}
\glsdefine{cobit}{COBIT}{Control Objectives for Information and Related Technologies}{COBIT -- Control Objectives for Information and Related Technologies, IT governance framework by ISACA distinguishing governance from management.}
\glsdefine{fair}{FAIR}{Factor Analysis of Information Risk}{FAIR -- Factor Analysis of Information Risk, a quantitative methodology for assessing and managing information risk.}
\glsdefine{isms}{ISMS}{Information Security Management System}{Information Security Management System -- the set of policies and procedures for systematically managing security (ISO 27001).}
\glsdefine{zt}{Zero Trust}{Zero Trust architecture}{Zero Trust -- security model based on never trust, always verify; access decisions based on identity and device posture rather than network location (NIST SP 800-207).}

% --- European Regulations and Legal Instruments ---
\glsdefine{gdpr}{GDPR}{General Data Protection Regulation}{General Data Protection Regulation (EU) 2016/679 -- EU regulation on the protection of personal data.}
\glsdefine{nis2}{NIS2}{NIS2 Directive}{NIS2 Directive (EU) 2022/2555 -- directive on measures for a high common level of cybersecurity.}
\glsdefine{aiact}{AI Act}{European AI Act}{European AI Act -- Regulation (EU) 2024/1689 on artificial intelligence, risk-based classification of AI systems.}
\glsdefine{cjeu}{CJEU}{Court of Justice of the European Union}{Court of Justice of the European Union -- the EU highest court, interpreting EU law.}
\glsdefine{ecthr}{ECtHR}{European Court of Human Rights}{European Court of Human Rights -- court enforcing the ECHR, based in Strasbourg.}
\glsdefine{edpb}{EDPB}{European Data Protection Board}{European Data Protection Board -- EU body ensuring consistent application of the GDPR.}
\glsdefine{edps}{EDPS}{European Data Protection Supervisor}{European Data Protection Supervisor -- independent EU authority overseeing the application of GDPR by EU institutions and bodies (distinct from the EDPB which coordinates national DPAs).}
\glsdefine{enisa}{ENISA}{European Union Agency for Cybersecurity}{ENISA -- European Union Agency for Cybersecurity.}
\glsdefine{dpo}{DPO}{Data Protection Officer}{Data Protection Officer -- person responsible for overseeing GDPR compliance within an organisation.}
\glsdefine{icradp}{ICRADP}{Independent Commission for Reasonable Adjustments of Digital Policies}{ICRADP -- Independent Commission for Reasonable Adjustments of Digital Policies (proposed in this report; French: CARPN).}

% --- Cross-border data transfers ---
\glsdefine{tia}{TIA}{Transfer Impact Assessment}{Transfer Impact Assessment -- documented assessment of whether a third country offers a level of personal-data protection essentially equivalent to EU law, required by CJEU Schrems II for transfers outside the EEA when relying on Standard Contractual Clauses.}
\glsdefine{scc}{SCC}{Standard Contractual Clauses}{Standard Contractual Clauses -- pre-approved European Commission contractual clauses for transfers of personal data outside the EEA; valid under Schrems II only when supplemented by a Transfer Impact Assessment and supplementary measures.}

% --- Sovereignty grid and technical reference vocabulary ---
\glsdefine{it}{IT}{Information Technology}{Information Technology -- the systems, software, networks and processes used to store, manipulate and transmit institutional information.}
\glsdefine{scim}{SCIM}{System for Cross-domain Identity Management}{System for Cross-domain Identity Management -- IETF standard (RFC 7643/7644) for automating the exchange of user identity information between identity domains.}
\glsdefine{byok}{BYOK}{Bring Your Own Key}{Bring Your Own Key -- a key-management arrangement in which the customer generates and provides the encryption keys used by a third-party service, retaining custody and rotation control.}
\glsdefine{hyok}{HYOK}{Hold Your Own Key}{Hold Your Own Key -- an arrangement stronger than BYOK in which the customer retains the keys in their own infrastructure (typically an HSM) and the service never has cleartext access; cryptographic operations require live customer participation.}
\glsdefine{rfp}{RFP}{Request for Proposal}{Request for Proposal -- a procurement document issued by an institution to solicit proposals from suppliers, specifying functional, technical, and contractual requirements.}
\glsdefine{iot}{IoT}{Internet of Things}{Internet of Things -- networked devices typically embedded in physical objects (sensors, controllers, lab instruments) with limited or no user interface and constrained security update mechanisms.}
\glsdefine{tls}{TLS}{Transport Layer Security}{Transport Layer Security -- IETF cryptographic protocol (RFC 8446 for TLS 1.3) providing confidentiality and integrity for network communications; the foundation of HTTPS.}
\glsdefine{opa}{OPA}{Open Policy Agent}{Open Policy Agent -- CNCF graduated open-source policy decision engine using the Rego language; widely adopted as a general-purpose authorisation and admission-control PDP.}
\glsdefine{pdp}{PDP}{Policy Decision Point}{Policy Decision Point -- the component that evaluates a request against the policy corpus and returns a decision; distinct from the Policy Enforcement Point (PEP) that applies the decision.}
\glsdefine{pep}{PEP}{Policy Enforcement Point}{Policy Enforcement Point -- the component that intercepts an access request, queries the PDP, and applies the returned decision (allow/deny plus obligations).}

% --- Research Networks ---
\glsdefine{csirt}{CSIRT}{Computer Security Incident Response Team}{Computer Security Incident Response Team -- team responsible for detecting, analysing, and responding to cybersecurity incidents.}
\glsdefine{eduroam}{eduroam}{Education Roaming}{eduroam -- Education Roaming, federated Wi-Fi service allowing users to connect at any participating institution worldwide using their home credentials.}
\glsdefine{edugain}{eduGAIN}{eduGAIN interfederation service}{eduGAIN -- global interfederation service connecting identity federations, enabling cross-institutional single sign-on.}
\glsdefine{geant}{GEANT}{GEANT pan-European research network}{GEANT -- pan-European backbone network connecting 40+ NRENs, serving 50 million users in 10,000+ institutions.}
\glsdefine{nren}{NREN}{National Research and Education Network}{National Research and Education Network -- a specialised ISP dedicated to supporting research and education (each country operates one).}
\glsdefine{first}{FIRST}{Forum of Incident Response and Security Teams}{FIRST -- Forum of Incident Response and Security Teams, global network of CSIRTs facilitating coordination across national and sectoral incident-response teams.}


% === Localization (key/value layer 0: generic defaults) ===
\usepackage{localization}
% ============================================================================
% localization-defaults.tex — LAYER 0 default bindings
% ----------------------------------------------------------------------------
% One \setloc per localization variable. This file is the authoritative registry
% of localization variables (49 distinct keys — INSTITUTION_NAME and
% ACADEMIC_GOVERNING_BODIES each support two anchor families and bind once here).
%
% Each default value is the variable's CANONICAL GENERIC ROLE PHRASE: the
% jurisdiction-neutral, institution-agnostic wording already used in the
% suite's current generic prose (the genericDecision of the matching row).
% Rendering a document with ONLY these defaults loaded therefore reproduces
% the present generic text UNCHANGED. An edition file (e.g. an EU or a
% Luxembourg edition) overrides any subset of these keys with \setloc after
% this file is \input.
%
% Keys are kebab-case. Requires localization.sty (\setloc / \loc).
% ============================================================================

% --- Subject institution & profile ------------------------------------------
% edition-name is the edition identity rendered prominently on every title
% page (\suiteditionline) and in \suitmetablock. The Layer-0 value makes an
% uninstantiated build self-identify loudly rather than pass for an edition.
\setloc{edition-name}{Generic reference edition (not instantiated)}
\setloc{institution-name}{the institution}
\setloc{institution-type}{a typical research-intensive institution}
\setloc{institution-macro-set}{the institution}
\setloc{staffing-profile}{an institution that sets excellence as a strategic objective}

% --- Institutional units & pilots -------------------------------------------
\setloc{cs-research-unit}{a computer-science department}
\setloc{security-research-centre}{a cybersecurity-focused interdisciplinary research centre}
\setloc{institution-hpc-facility}{the institution's HPC facility / HPC team}
\setloc{pilot-entities}{the institution's HPC facility / team}

% --- Institutional governance roles -----------------------------------------
\setloc{executive-leadership}{the institution's central executive leadership}
\setloc{governing-board}{the institution's governing board}
\setloc{academic-governing-bodies}{the academic deliberative councils}
\setloc{academic-endorsement-path}{the relevant academic deliberative councils}
\setloc{institution-faculty-structure}{typically organised into faculties, each with departments of teaching and research, and interdisciplinary centres}
\setloc{staff-representation-body}{the institution's staff-representation bodies}
% Approving academic units for the internal-approval byline on the dossier and
% the rectorate summary (the faculty/department domains that signed off the
% institution's edition). Generic default is the neutral role phrase.
\setloc{approving-academic-units}{the approving academic units of \loc{institution-name}}

% Suite version (generic, maintainer-managed; bumped by tools/bump-version.sh).
% ============================================================================
% suite-version.tex --- SINGLE SOURCE OF TRUTH for the generic SUIT suite
% version (maintainer-managed). Common to all four generic documents.
% ----------------------------------------------------------------------------
% Semantic version X.Y.Z. Bumped ONLY by tools/bump-version.sh:
%   patch (default) -> Z+1   (one bump per release, regardless of change count)
%   minor           -> Y+1, Z=0   (maintainer, on request)
%   major           -> X+1, Y=0, Z=0   (maintainer, on request)
% Do not edit by hand other than through the bump tool; the line below is
% matched verbatim by the tool's parser.
% ============================================================================
\setloc{suite-version}{1.0.0}

% Per-document edition versions: EMPTY in the generic build (the suite version
% is shown instead); an edition fills, in its set.tex, the documents it
% publishes. Keyed by document id (derived from \jobname in \suitmetablock).
\setloc{policy-EN-edition-version}{}
\setloc{policy-summary-EN-edition-version}{}
\setloc{solution-EN-edition-version}{}
\setloc{solution-summary-EN-edition-version}{}

% Document metadata for the unified \suitmetablock title-page block. Each
% edition overrides these for its own visibility/approval facts.
\setloc{doc-visibility}{INTERNAL}
\setloc{doc-approval-authority}{\loc{approving-academic-units}}
\setloc{doc-approval-date}{\textit{(approval date to be set)}}

% --- Research & education network and its services --------------------------
\setloc{nren-name}{the national research and education network}
\setloc{nren-csirt}{the NREN CSIRT}
\setloc{national-private-sector-csirt}{a separate CSIRT serving the private sector, communes and non-governmental entities}
\setloc{financial-sector-supervisor}{a financial-sector supervisor}
\setloc{research-ecosystem-partners}{a handful of national research institutes}
\setloc{sdmz-status-lead}{In many national ecosystems, the}

% --- Companion openness-status table: one marker per element ('\warn' = to be
%     assessed in the generic build; an edition binds the institution's verdict) ---
\setloc{status-table-subtitle}{current status}
\setloc{status-byod}{\warn}
\setloc{status-eduroam}{\warn}
\setloc{status-mobility}{\warn}
\setloc{status-internet}{\warn}
\setloc{status-cloud-ai}{\warn}
\setloc{status-sci-sources}{\warn}
\setloc{status-local-admin}{\warn}
\setloc{status-tool-choice}{\warn}
\setloc{status-environments}{\warn}
\setloc{status-data-platforms}{\warn}
\setloc{status-code-repos}{\warn}
\setloc{status-open-access}{\warn}
\setloc{status-labs}{\warn}
\setloc{status-cyber-ranges}{\warn}
\setloc{status-api-access}{\warn}

% --- Abstract genesis + ICRADP adjustments-committee precedent ---
\setloc{abstract-genesis-institution}{a research-intensive university}
\setloc{existing-adjustments-committee-precedent}{Many institutions already operate a \keyword{reasonable-adjustments committee} (or equivalent)}
\setloc{icradp-committee-precedent-lead}{Where the institution already operates a reasonable-adjustments mechanism (supporting students and staff), the \gls{icradp} extends that principle to digital policies; where none exists, it establishes the principle directly.}
\setloc{cite-adjustments-committee}{}
\setloc{cite-equal-treatment-law}{}
\setloc{funding-instrument}{multi-year planning and funding cycle}
\setloc{peer-nren-roster}{peer NRENs in comparable jurisdictions, each operating a national backbone, a CSIRT, and federated identity/storage/IaaS services}
\setloc{nren-service-catalogue}{the NREN-provided federated-service catalogue (federated Wi-Fi, trusted certificate service, performance monitoring, managed research-data transfer)}
\setloc{nren-research-cloud}{a national/regional research cloud or NREN-provided cloud service, where one exists}
\setloc{research-backbone-scale}{the regional research backbone (membership and backbone capacity per region)}
\setloc{identity-federation-profile}{the academic identity interfederation}

% --- High-performance computing ---------------------------------------------
\setloc{national-hpc-facility}{a national/EuroHPC supercomputer}

% --- Deployment topology & data custody -------------------------------------
\setloc{deployment-custody-model}{the institution operates, or procures under equivalent custody guarantees, a deployment of its core components}
\setloc{data-custody-sovereignty-profile}{a custody arrangement that minimises the applicable disclosure-regime exposure}

% --- Regulatory regimes (supranational anchors as default profile) ----------
\setloc{regulatory-regime}{binding law in the EU and Council of Europe member states (including the institution's jurisdiction)}
\setloc{cybersecurity-baseline-regime}{the applicable supranational and national legal framework}
\setloc{data-protection-regime}{the institution's applicable data-protection regime}
\setloc{eid-trust-services-regime}{the institution's applicable trust-services regime}
\setloc{ai-regulation}{the adopter's applicable AI-governance regime}
\setloc{cross-border-transfer-regime}{the applicable cross-border data-transfer regime}
\setloc{essential-entity-classification-basis}{, the public-administration entry of NIS2 Annex~I being the basis in the worked example}
\setloc{nis2-sanction-ceiling}{up to EUR\,10\,M or 2\,\% of annual turnover}
% --- NIS2 essential-entity status (conditional-fact pattern; an edition asserts the established fact via its set) ---
\setloc{essential-entity-modality}{may be classified}
\setloc{nis2-status-conditional-tail}{, depending on the national transposition}
\setloc{cer-designation-status}{A given institution may or may not be designated as a CER critical entity}
% --- Edition-routed citations: empty in the generic build; bound to \cite{} markers in an edition set ---
\setloc{cite-ilr-essential}{}
\setloc{cite-nis2-transposition}{}
\setloc{cite-notification-platform}{}
\setloc{cite-national-risk-method}{}
\setloc{cite-public-sector-isp}{}
\setloc{cite-cyber-strategy}{}
\setloc{cite-nren-csirt}{}
\setloc{cite-cer-transposition}{}
\setloc{cite-gov-csirt}{}
\setloc{cite-spoc}{}
\setloc{cite-infosec-agency}{}
\setloc{cite-circl}{}
\setloc{cite-cssf}{}
\setloc{cite-restena}{}

% --- National transposition instruments -------------------------------------
\setloc{nis2-national-transposition}{the national NIS2 transposition law}
\setloc{cer-national-transposition}{the national CER transposition law}
\setloc{national-equal-treatment-law}{the national transposition of Directive 2000/78/EC into labour law}
\setloc{national-public-sector-isp}{a national public-sector information-security policy}
\setloc{national-cyber-strategy}{the national cybersecurity strategy}
\setloc{national-risk-method}{a national risk-analysis method}
\setloc{national-cyber-commissariat-founding-law}{the national founding statute of the cyber-resilience commissariat}

% --- National authorities & bodies (role-based) -----------------------------
\setloc{national-dpa}{the competent data-protection authority}
\setloc{national-gov-csirt}{the national/governmental CSIRT}
\setloc{national-infosec-agency}{the national information-security agency}
\setloc{nis2-supervisory-authority}{the national NIS2 supervisory authority}
\setloc{national-spoc-authority}{the national single point of contact}
\setloc{national-cyber-competence-centre}{a national cybersecurity competence centre}
\setloc{national-notification-platform}{the national NIS2 incident-notification platform}
\setloc{national-legislature}{the national legislature}
\setloc{national-executive-department}{the national executive department responsible for the protection commissariat}

% --- Operational contacts & market exemplars --------------------------------
\setloc{internal-emergency-incident-contact}{the institution's internal emergency / incident-reporting contact}
\setloc{incumbent-identity-suite-vendor}{a pure-commercial identity provider}

% --- ICRADP administrative-review service levels ----------------------------
% Adjustment-commission decision deadline and emergency provisional-grant
% window. Default to the worked figures; an edition may rebind them to its own
% administrative-review service levels.
\setloc{icradp-decision-deadline}{15 working days}
\setloc{icradp-emergency-deadline}{48 hours}

% --- Assurance / review recency window (single-source magic number) ---------
% The audit/pen-test/drill recency window used by the Chapter 4 self-assessment.
% Centralised so a re-tuned interval is one edit and an edition may rebind it.
\setloc{assurance-recency-window}{twenty-four months}

\endinput


% ============================================================================
% CUSTOM COMMANDS (aligned with the rest of the suite)
% ============================================================================
\newcommand{\keyword}[1]{\textbf{#1}}
\newcommand{\institution}[1]{\textsc{#1}}
\newcommand{\framework}[1]{\textsf{#1}}

% Five-layer evidence markers (legal > standards > sector > empirical > institutional)
\newcommand{\evlegal}{\textcolor{BrickRed}{\textbf{[Legal]}}\space}
\newcommand{\evstd}{\textcolor{NavyBlue}{\textbf{[Standards]}}\space}
\newcommand{\evsector}{\textcolor{Plum}{\textbf{[Sector]}}\space}
\newcommand{\evempirical}{\textcolor{OliveGreen}{\textbf{[Empirical]}}\space}
\newcommand{\evinst}{\textcolor{Slate}{\textbf{[Institutional]}}\space}

% Cross-reference to a chapter of the Technical Reference being defended
\newcommand{\trref}[1]{the Technical Reference (\emph{#1})}

% ============================================================================
% CUSTOM TCOLORBOX ENVIRONMENTS (suite question/answer convention)
% ============================================================================
\newtcolorbox{questionbox}[1][]{
    colback=red!3!white,
    colframe=red!60!black,
    fonttitle=\bfseries,
    title={#1},
    breakable,
    sharp corners=south,
    boxrule=0.6pt,
    left=4pt, right=4pt
}

\newtcolorbox{answerbox}[1][]{
    colback=OliveGreen!3!white,
    colframe=OliveGreen!60!black,
    fonttitle=\bfseries,
    title={#1},
    breakable,
    sharp corners=south,
    boxrule=0.6pt,
    left=4pt, right=4pt
}

\newtcolorbox{clarificationbox}[1][]{
    colback=NavyBlue!3!white,
    colframe=NavyBlue!60!black,
    fonttitle=\bfseries,
    title={#1},
    breakable,
    sharp corners=south,
    boxrule=0.6pt,
    left=4pt, right=4pt
}

% ============================================================================
% SECTION NUMBERING: skip chapter level (we use \part > \section directly)
% ============================================================================
\renewcommand{\thesection}{\arabic{section}}
\renewcommand{\thesubsection}{\thesection.\arabic{subsection}}
\renewcommand{\thesubsubsection}{\thesubsection.\arabic{subsubsection}}

% === Prevent vertical stretching ===
\raggedbottom

% === Part title on right page, content starts immediately on next page ===
\makeatletter
\renewcommand{\@endpart}{\vfil\newpage}
\makeatother

% ============================================================================
\begin{document}

% === Title page ===
\hypersetup{pageanchor=false}
\begin{titlepage}
\centering
\vspace*{0.5cm}

{\suittitlelogo}

\vspace{0.8cm}

{\Huge\bfseries Sustainable University IT}

\vspace{0.5cm}
{\LARGE\itshape Technical and Operational\\Objections and Responses}

\vspace{1.2cm}
{\large An anticipatory argumentation module defending\\
the reference architecture, the sovereignty grid,\\
and the brownfield migration playbook}

\vspace{2cm}

\begin{tcolorbox}[
    colback=NavyBlue!8!white,
    colframe=NavyBlue!85!black,
    boxrule=1.2pt,
    arc=3pt,
    left=12pt, right=12pt, top=8pt, bottom=8pt,
    width=0.88\textwidth,
    halign=center
]
{\large\bfseries Companion to\\
``\suitsolutiontitle\ --- \suitsolutionsubtitlemain\\
\suitsolutionsubtitlerest''}
\end{tcolorbox}

\vspace{1cm}
\suitmetablock

\vspace{0.9cm}

\begin{minipage}{0.85\textwidth}
\centering
{\large\bfseries SUIT}\\[3pt]
{\normalsize\itshape International Working Group\\on Sustainable University IT in the AI Era}\\[4pt]
{\small\itshape (in formation)}
\end{minipage}

\vspace{0.9cm}

\begin{minipage}{0.85\textwidth}
\centering\small\itshape
This module addresses good-faith technical and operational objections
to the Technical Reference. It is intended for chief technology officers,
information-security officers, network and platform engineers, procurement
officers, and the governance bodies that mandate them. Every response is
constructed on five independent evidence layers --- legal, standards,
sector, empirical, and institutional --- and remains generic across
institutions and jurisdictions.
\end{minipage}

\vspace{0.6cm}

\begin{minipage}{0.85\textwidth}
\small\noindent\textbf{Note on the process:}\\
\textit{This document is an always living document. It is provided for
constructive feedback and revision in order to improve its quality. The
engineering of this document was partially supported by digital tools, some
of which incorporate Artificial Intelligence features.}\\[3pt]
\textcolor{Crimson}{\textit{This document contains a portion of imprecisions,
missing content, errors or possibly hallucinations. We believe that a
supportive and cooperative approach by the readers will make its content
quality increase over time.}}
\end{minipage}

\vspace{1.5cm}
\end{titlepage}

% === Frontmatter ===
\frontmatter
\hypersetup{pageanchor=true}

\chapter*{How to read this module}
\addcontentsline{toc}{chapter}{How to read this module}

This module anticipates the \keyword{technical and operational}
objections most likely to be raised against \trref{Architecture and
Migration}. It deliberately excludes political, budgetary-in-principle,
and academic-freedom objections, which are treated in the policy-layer
argumentation; here the question is always the same: \emph{does the
proposed engineering hold up under adversarial scrutiny?}

Each objection is stated in a \textcolor{red!60!black}{\textbf{question
box}} in the words a sceptical engineer or architect would actually use.
The response is given in a \textcolor{OliveGreen!60!black}{\textbf{answer
box}} that follows a strict \keyword{concession-then-rebuttal} pattern:
it first concedes the legitimate core of the objection, then rebuts the
overreach. Where a point needs disambiguation rather than refutation, an
\textcolor{NavyBlue!60!black}{\textbf{clarification box}} is used instead.

\paragraph{The five-layer evidence methodology.}
No response rests on assertion. Each rebuttal is built from up to five
\emph{independent} evidence layers, marked inline so the reader can audit
the chain of support:

\begin{description}[leftmargin=2.2cm,style=nextline]
  \item[\evlegal] binding legal and regulatory instruments of the
        adopter's jurisdiction (\institution{GDPR},
        \institution{NIS2} and sectoral law in the worked EU
        instantiation).
  \item[\evstd] open technical standards and normative frameworks
        (\institution{IETF}, \institution{OASIS}, \institution{NIST},
        \institution{ISO}).
  \item[\evsector] higher-education and research-sector practice,
        toolkits, and sector bodies.
  \item[\evempirical] measured outcomes --- breach data, incident
        reports, performance studies, cost figures.
  \item[\evinst] named institutional precedents operating the relevant
        pattern in production.
\end{description}

A response is considered sound when at least \emph{three} of these layers
converge independently --- the same minimum-three-layer rule applied
throughout the suite.

\tableofcontents

% ============================================================================
\mainmatter
% ============================================================================

\part{Technical and Operational Objections}

% ----------------------------------------------------------------------------
\section{Architecture complexity}
\label{sec:arg-complexity}

\begin{questionbox}[O1.1 --- ``The nine-layer architecture is over-engineered for a university.'']
The reference architecture decomposes the estate into nine layers ---
identity, network, compute, storage, observability, cryptography,
orchestration, policy, endpoint --- each with its own interface contract.
A university is not a hyperscaler. This is enterprise-architecture theatre
that no campus IT team can realistically operate; the layering adds
cognitive load and failure modes without adding value.
\end{questionbox}

\begin{answerbox}[Response O1.1]
\textbf{Concession.} The objection is right that \emph{a large number of
moving parts is itself a risk}. Architectural complexity that is not
matched by operational capacity degrades reliability rather than improving
it; this is the central lesson of dependable-systems engineering, where
each added component is an added fault source unless it is contained by a
clear interface \citep{avizienis2004}.

\medskip
\textbf{Rebuttal.} The objection mistakes \emph{decomposition} for
\emph{addition}. The nine layers are not nine new systems; they are nine
\emph{names for functions that already exist} in every university estate ---
identity, network, storage, and so on are present whether or not they are
drawn as layers. The architecture does not add them; it makes their
\keyword{interfaces explicit} so that complexity becomes \emph{bounded and
auditable} instead of tangled.

\begin{itemize}[leftmargin=1.4em,itemsep=2pt]
  \item \evstd Layered decomposition behind stable interfaces is the
        canonical method for managing, not multiplying, complexity in
        secure systems engineering \citep{nist-sp800-160v2,avizienis2004}.
        \institution{Zero Trust} architecture is itself defined as a set of
        separated planes (policy decision, policy enforcement, data) for
        exactly this reason \citep{nist-zt}.
  \item \evsector The sector's own governance frameworks treat
        function-by-function decomposition as the prerequisite for
        manageability, not its enemy \citep{cobit,dehaes2015}.
  \item \evinst Research-intensive institutions already operate this decomposition in
        production --- separated identity, network, and storage planes are
        visible, for example, at \institution{MIT} \citep{mit-ist,mit-kerberos},
        \institution{Stanford} \citep{stanford-uit,stanford-sunet}, and
        \institution{University of Michigan} \citep{umich-its,umich-core-network}.
        None is a hyperscaler; all operate a layered decomposition along the
        same functional planes, whatever local labels they use.
\end{itemize}

\textbf{Bottom line.} The architecture does not impose nine systems on a
campus that had one. It gives names and contracts to the functions a campus
already runs, so the irreducible complexity is \emph{governed} rather than
\emph{hidden}.
\end{answerbox}

\begin{clarificationbox}[Clarification --- complexity is instantiated to scale]
The nine-layer \emph{model} is invariant; its \emph{instantiation} is not.
A teaching-led institution may realise several layers as a single managed
service or a hosted offering, collapsing operational surface while keeping
the interface contracts intact (\trref{Sizing reference points}). The model
is the vocabulary; the deployment is sized to the institution.
\end{clarificationbox}

\begin{questionbox}[O1.2 --- ``Identity and policy become institution-wide single points of failure, with no availability story --- a commercial SLA beats this.'']
Per-session policy decision routes every activation through the
identity and policy planes. If either is down, nobody works. A
managed commercial identity service at least carries a contractual
SLA for that risk; the Reference proposes to self-operate the most
critical planes of the institution.
\end{questionbox}

\begin{answerbox}[Response O1.2]
\textbf{Concession.} The mediation-point risk is real, and it is
\emph{created by the architecture itself}: making activation
explicit concentrates what a flat estate leaves diffuse. Any
deployment that treated the identity provider or the policy
decision point as ordinary single-instance services would deserve
this objection in full.

\medskip
\textbf{Rebuttal.} The Reference engineers the risk rather than
insuring against it (\trref{Architectural principles}, availability
and failure semantics of the mediation planes): each critical plane
runs redundantly with a declared availability target measured
through the observability layer; policy enforcement points cache
decisions for the activation-token lifetime, so a policy-plane
outage blocks \emph{new} activations without invalidating
established sessions; and each zone class declares its fail mode
(fail-closed for administrative and sensitive zones, fail-degraded
with audit where connectivity availability prevails).

\begin{itemize}[leftmargin=1.4em,itemsep=2pt]
  \item \evstd Zero-trust guidance itself requires exactly this
        engineering: NIST SP~800-207 treats PDP/PEP availability
        and decision caching as design obligations of the model
        \citep{nist-zt}.
  \item \evstd The reference components are built for clustered
        operation --- replicated identity providers, policy engines
        evaluating from locally distributed bundles
        \citep{opa-project} --- so redundancy is configuration, not
        custom engineering.
  \item \evinst The failure semantics are testable, not asserted:
        availability targets are measured by the observability
        layer, and the fail-mode declarations are part of each
        zone's published security profile, exercised through the
        conformance suite (\trref{Conformance and replaceability
        tests}).
\end{itemize}

\textbf{Bottom line.} An SLA is a remedy after the fact; declared,
rehearsed failure semantics are a design. The architecture does not
leave the mediation planes' availability to contract --- it
engineers it, measures it, and makes outage behaviour a
per-zone-class policy choice.
\end{answerbox}

% ----------------------------------------------------------------------------
\section{Open-source coverage gaps}
\label{sec:arg-oss-gaps}

\begin{questionbox}[O2.1 --- ``Open-source MDM cannot manage iOS and Android, so the open-foundations claim is false.'']
The whole document rests on a reference instantiation built on open
foundations. But there is no credible open-source mobile-device-management
stack for \institution{iOS} and \institution{Android} at scale. The moment
a university has a real fleet of phones and tablets, it must buy a
proprietary suite. The ``open foundations'' claim collapses at the endpoint.
\end{questionbox}

\begin{answerbox}[Response O2.1]
\textbf{Concession.} This is the strongest technical objection in the
document, and it is conceded \emph{in full on the facts}. There is no mature
open-source MDM that matches a commercial suite for supervised
\institution{iOS} and managed \institution{Android} at fleet scale. Open
tooling for Apple platforms is real and production-grade for enrolment and
software management \citep{micromdm,munki}, and \institution{osquery}
provides open posture telemetry across operating systems
\citep{osquery}; but a turnkey, vendor-supported mobile MDM on fully open
foundations does not exist today. The Technical Reference says exactly this,
in the open, rather than hiding it.

\medskip
\textbf{Rebuttal.} The objection then over-reaches: it treats one bounded
gap as the collapse of the whole. The open-foundations claim is a claim
about \keyword{replaceability through contracts}, not about purism. The
architecture's test is not ``is every component open?'' but ``can every
component be replaced without re-architecting?'' --- and a proprietary MDM
that exports configuration profiles to a documented format and feeds posture
to the policy plane through a standard interface \emph{passes that test}.

\begin{itemize}[leftmargin=1.4em,itemsep=2pt]
  \item \evstd The endpoint contract is defined by open interfaces, not by
        the product behind them: posture and identity flow over
        \institution{SCIM} \citep{rfc-7643-scim-core}, \institution{OIDC}
        \citep{oidc-core-1.0}, and standard telemetry --- so a proprietary
        suite is a conforming \emph{instantiation}, not an exception.
  \item \evsector The sector already procures exactly this way: the
        \institution{HECVAT} toolkit exists precisely because universities
        buy assessed proprietary endpoint products under a common
        contract \citep{hecvat}. Commercial suites such as
        \citep{jamf,intune} are named in the reference as legitimate
        instantiations for the iOS/Android portion of the estate.
  \item \evempirical Endpoint compromise --- not MDM ideology --- is where
        the measured risk lives: ransomware and credential-driven breaches
        dominate education-sector incident data
        \citep{sophos-ransomware-edu2025,verizon-dbir2024}. The architecture
        prioritises \emph{posture and replaceability} over provenance
        precisely because that is what the breach data rewards.
  \item \evinst Mixed open/proprietary endpoint estates are common at
        research-intensive institutions, which run a dedicated Apple-platform suite
        alongside open configuration management and \institution{osquery}
        for Linux and servers (\institution{MIT}, \institution{Stanford},
        \institution{University of Michigan}
        \citep{mit-ist,stanford-uit,umich-safecomputing}).
\end{itemize}

\textbf{Bottom line.} The iOS/Android MDM gap is real and acknowledged. It
does not falsify the open-foundations claim, because that claim is about
\keyword{contract-bounded replaceability}, and a proprietary MDM that honours
the endpoint interface contract is a conforming instantiation, not a
counter-example.
\end{answerbox}

\begin{questionbox}[O2.2 --- ``Open source is itself the supply-chain risk: unmaintained dependencies, no vendor CVE-response SLA, nobody to sue.'']
The open-foundations default trades a contractual security
relationship for a volunteer one. Open components ship unmaintained
transitive dependencies, carry no committed CVE-response time, and
offer no counterparty bearing liability when something breaks.
\end{questionbox}

\begin{answerbox}[Response O2.2]
\textbf{Concession.} The failure mode is real: abandonware exists,
and an open licence confers no right to a fix. A deployment that
consumed anonymous code with no support arrangement, no provenance
control, and no maintenance assessment would deserve the objection
in full.

\medskip
\textbf{Rebuttal.} The architecture does not consume anonymous
code; it procures open components under the same discipline as any
supplier.

\begin{itemize}[leftmargin=1.4em,itemsep=2pt]
  \item \evsector Where a committed CVE-response time is required,
        the institution buys a commercial support subscription for
        the open component --- procured and assessed like any
        vendor relationship \citep{hecvat}; the counterparty
        exists, and what remains open is the exit.
  \item \evstd Supply-chain provenance is a named control, not a
        hope: SLSA provenance for built artefacts
        \citep{slsa-spec}, with the supply-chain-security ownership
        established for NIS2 Art.~21(2)(d) by the policy module's
        Q8.5.
  \item \evinst The sovereignty grid's \emph{Operational} dimension
        scores maintenance reality --- skills, community,
        certification paths --- per layer \emph{before} adoption
        (\trref{The sovereignty grid}); a component nobody can
        sustain scores itself out.
  \item \evstd ``Nobody to sue'' cuts both ways: a contractual
        remedy does not produce a patch, while source access and
        the replaceability conditions (O5.1) let the institution
        patch, fork, or replace --- bounding the blast radius of
        any single project's abandonment.
\end{itemize}

\textbf{Bottom line.} Supply-chain risk attaches to unmanaged
consumption, not to licence type. The architecture pairs open
components with support contracts where SLAs matter, provenance
controls where builds matter, and an operational-sustainability
score before adoption --- with replaceability as the final
backstop.
\end{answerbox}

% ----------------------------------------------------------------------------
\section{Sizing realism}
\label{sec:arg-sizing}

\begin{questionbox}[O3.1 --- ``The sizing reference points are fantasy; a small university cannot run this.'']
The architecture may suit a flagship research university with hundreds of
network engineers, but it is unaffordable and unstaffable for a small
teaching-led institution. Prescribing a nine-layer estate to a college with
a five-person IT team is sizing fantasy.
\end{questionbox}

\begin{answerbox}[Response O3.1]
\textbf{Concession.} Conceded that \emph{a one-size instantiation would be
fantasy}. A small institution genuinely cannot operate the same realisation
as a multi-petabyte research university, and any document that demanded it
would be unserious about operations.

\medskip
\textbf{Rebuttal.} The architecture does not demand a single realisation. It
distinguishes the \keyword{minimal necessary functional set} --- which is
size-invariant --- from its \keyword{instantiation} --- which scales. The
functions an institution must perform (authenticate subjects, segment the
network, classify and protect data, retain auditable logs) do not disappear
at small scale; they are simply realised more compactly, often as managed or
hosted services. This is a guardrail \emph{against} overclaim, not an
instance of it.

\begin{itemize}[leftmargin=1.4em,itemsep=2pt]
  \item \evlegal The size-invariant core is anchored in law ---
        \loc{data-protection-regime} and
        \loc{cybersecurity-baseline-regime}, worked here through the
        EU instruments: \institution{GDPR}
        imposes accountability, security of processing, and breach
        obligations with \emph{no de-minimis threshold} ---
        a five-person institution owes the same duties as a flagship
        \citep{gdpr}. (The one narrow size-based relief, the
        records-of-processing derogation for organisations below
        250~persons under Article~30(5), is itself conditional and does not
        touch the security, accountability, or breach duties at issue here.)
        \institution{NIS2} likewise scales nothing down: once an
        entity is in scope (by sector-and-size, or by a
        size-independent designation route), its duties apply in
        full regardless of headcount \citep{nis2}.
  \item \evstd The minimal set maps to baseline controls that are stated
        independently of scale \citep{nist-csf,iso27001,nist-zt}; the
        controls are mandatory, the deployment topology is sized.
  \item \evsector \institution{EDUCAUSE} governance guidance frames IT
        provision as a function of mission and risk, explicitly supporting
        shared-service and consortial realisations for smaller institutions
        \citep{educause-governance,educause-isg}.
  \item \evinst National research and education networks
        (e.g.\ \institution{GÉANT}, \institution{RENATER}, \institution{RESTENA}
        in Europe) exist precisely so that small institutions consume identity,
        connectivity, and security as shared services rather than building
        them alone \citep{geant,renater,restena,eduroam,edugain}.
\end{itemize}

\textbf{Bottom line.} The minimal functional set is size-invariant because
the \emph{obligations} are; the instantiation scales down to managed and
shared services. The sizing chapter prescribes envelopes, not a single
mandatory plant.
\end{answerbox}

% ----------------------------------------------------------------------------
\section{Brownfield migration and coexistence}
\label{sec:arg-migration}

\begin{questionbox}[O4.1 --- ``This is a greenfield design; no real university can rip-and-replace its estate.'']
The wave-based playbook reads like a clean-sheet build. Real universities
have decades of accreted systems, in-flight contracts, and departments that
will not be coordinated. A migration that assumes you can stop the world and
rebuild is dead on arrival.
\end{questionbox}

\begin{answerbox}[Response O4.1]
\textbf{Concession.} Conceded that \emph{a rip-and-replace migration would
fail}. Estates are heterogeneous, contracts are mid-term, and federated
governance means no central team can mandate a flag day. Any playbook that
ignored this would be unusable.

\medskip
\textbf{Rebuttal.} The objection describes a design the Technical Reference
explicitly rejects. The playbook is \keyword{brownfield by construction}: it
is wave-based precisely so that each step is independently valuable,
reversible, and \emph{coexists} with incumbent systems. A dedicated
coexistence-patterns library and a reversibility playbook exist for exactly
the heterogeneity the objection raises (\trref{Coexistence patterns library};
\trref{Reversibility playbook}).

\begin{itemize}[leftmargin=1.4em,itemsep=2pt]
  \item \evstd Incremental adoption behind stable interfaces is the
        standard path for cyber-resilient transformation
        \citep{nist-sp800-160v2}; \institution{Zero Trust} migration is
        itself specified as incremental and coexistence-friendly, never as
        a flag day \citep{nist-zt}.
  \item \evsector Sector governance treats migration as portfolio change
        managed against mission risk, not as a single cutover
        \citep{cobit,weill-ross,educause-governance}.
  \item \evempirical The waves are ordered to address the most-exploited
        attack patterns first: identity, segmentation, and observability
        lead because credential abuse dominates measured breach data
        \citep{verizon-dbir2024,ibm-cost-breach-2024} and segmentation
        interrupts its in-network progression \citep{nist-sp800-215}.
        The ordering is
        an evidence-informed prioritisation of the measured exposure, so
        early waves buy down the most-exploited risk before any large
        investment.
  \item \evinst Coexistence is the lived reality at research-intensive
        institutions such as \institution{MIT} and
        \institution{University of Michigan}, where federated
        identity overlays a heterogeneous estate \citep{mit-kerberos,umich-its},
        and \institution{Science DMZ} deployments coexist with the general
        campus network rather than replacing it
        \citep{science-dmz-2013,esnet-cases}.
\end{itemize}

\textbf{Bottom line.} The playbook is the opposite of greenfield: each wave
is independently valuable and reversible, and coexistence with incumbents is
a first-class design goal, not an afterthought.
\end{answerbox}

% ----------------------------------------------------------------------------
\section{Vendor lock-in versus replaceability}
\label{sec:arg-lockin}

\begin{questionbox}[O5.1 --- ``Standardising on a stack just swaps one lock-in for another.'']
You criticise proprietary lock-in, then prescribe a specific open stack ---
a particular identity provider, a particular orchestrator, a particular
storage system. That is still lock-in; you have only changed the vendor. The
``replaceability'' claim is rhetorical.
\end{questionbox}

\begin{answerbox}[Response O5.1]
\textbf{Concession.} Conceded that \emph{naming open products does not by
itself defeat lock-in}. An open component can still trap an institution
through bespoke configuration, undocumented data formats, or unique
operational dependencies. Open licensing is necessary but not sufficient for
replaceability; status-quo bias makes any incumbent --- open or not ---
sticky \citep{samuelson-zeckhauser}.

\medskip
\textbf{Rebuttal.} The architecture does not claim replaceability \emph{from}
choosing open products. It claims replaceability \emph{from} four explicit
conditions enforced by a conformance suite: standard interfaces, open
state and configuration formats, documented exit-data export, and no
hidden dependencies. The
named open instantiation is an \emph{example} that satisfies these
conditions; a proprietary component that satisfies the same conditions is
equally acceptable. Lock-in is defeated by the \keyword{contract}, not by the
\keyword{logo}.

\begin{itemize}[leftmargin=1.4em,itemsep=2pt]
  \item \evstd Every layer's interface is pinned to an open standard ---
        \institution{OIDC}/\institution{SAML}/\institution{SCIM} for identity
        \citep{oidc-core-1.0,saml-core-2.0,rfc-7643-scim-core},
        \institution{KMIP}/\institution{PKCS\#11} for key management
        \citep{kmip-oasis,pkcs11-oasis}, \institution{TLS}~1.3 and
        \institution{ACME} on the wire \citep{rfc-8446-tls13,rfc-8555-acme},
        \institution{OpenTelemetry} and \institution{IPFIX} for telemetry
        \citep{opentelemetry-spec,rfc-7011-ipfix}, and
        \institution{OCI} image formats for workloads \citep{oci-image-spec}.
        Replaceability is the property of conforming to these, regardless of
        the product.
  \item \evlegal Replaceability is also a compliance asset: \institution{GDPR}
        portability and the ability to change processors presuppose that data
        can leave a system in a usable form \citep{gdpr}; standard-contractual
        and processor-change mechanisms assume exit is possible
        \citep{eu-sccs2021}.
  \item \evsector \institution{HECVAT} institutionalises exactly this
        vendor-neutral, contract-based assessment for the sector
        \citep{hecvat}.
  \item \evinst Multi-vendor, standards-pinned identity federations
        (\institution{eduGAIN} and \institution{eduroam} globally, with regional
        federations such as \institution{InCommon})
        demonstrate that replaceability through interfaces is an operating
        reality, not a slogan \citep{edugain,incommon,eduroam}.
\end{itemize}

\textbf{Bottom line.} The defence against lock-in is the four replaceability
conditions and the conformance suite that tests them --- not allegiance to
any product, open or proprietary.
\end{answerbox}

% ----------------------------------------------------------------------------
\section{Conformance-test rigour}
\label{sec:arg-conformance}

\begin{questionbox}[O6.1 --- ``The conformance tests are a checklist, not a real guarantee.'']
``Replaceability conditions'' and a ``conformance day'' sound reassuring, but
a checklist is not proof. Without an enforced, repeatable, and independently
verifiable test, conformance is a self-graded exercise that every vendor will
claim to pass.
\end{questionbox}

\begin{answerbox}[Response O6.1]
\textbf{Concession.} Conceded that \emph{an unenforced checklist is
worthless}. A conformance claim that cannot be re-run by a third party, on
representative data, is marketing. Self-attestation without evidence is a
known failure mode in procurement.

\medskip
\textbf{Rebuttal.} The conformance suite is constructed as an
\keyword{executable, evidence-producing test}, not a questionnaire. It
defines categories of test --- interface fidelity, data portability, exit
rehearsal, and decision auditability --- each of which yields an artefact: a
captured export consumed by an alternative, a replay against a documented
schema, a recorded migration on a representative subset
(\trref{Conformance and replaceability tests}). The test is passed by
\emph{producing the artefact}, not by asserting compliance. Where
sound practice makes a nominal artefact unobtainable ---
non-exportable HSM trust-anchor keys being the canonical case ---
the layer contract declares the substitute artefact (a key-rotation
and re-anchoring rehearsal under dual control), so the gate is
passable without weakening key custody.

\begin{itemize}[leftmargin=1.4em,itemsep=2pt]
  \item \evstd Conformance-to-interface as the unit of verification is the
        standard discipline: the open interfaces themselves
        (\institution{SCIM}, \institution{KMIP}, \institution{OCI},
        \institution{OpenTelemetry}) carry normative conformance criteria
        that a test can mechanically check
        \citep{rfc-7643-scim-core,kmip-oasis,oci-image-spec,opentelemetry-spec}.
  \item \evsector The sector already operates standardised, repeatable
        third-party assessment through \institution{HECVAT}, and research
        cybersecurity centres publish reusable assessment methods
        \citep{hecvat,trustedci}.
  \item \evempirical Tests are weighted toward the data: control categories
        track the attack patterns that actually cause breaches --- exfiltration
        paths, credential reuse, and unpatched transfer software ---
        as reported in incident data
        \citep{verizon-dbir2024,moveit-breach}.
  \item \evinst The artefact-based discipline is already routine in the
        sector, if for a different property: \institution{perfSONAR} and
        \institution{Science DMZ} deployments verify end-to-end behaviour
        against measured baselines rather than vendor claims
        \citep{perfsonar,science-dmz-2013,sciencedmz-perf}. This establishes
        that re-runnable, evidence-producing conformance is an operating
        norm on the research network, the same discipline the replaceability
        suite applies to exit and portability.
\end{itemize}

\textbf{Bottom line.} Conformance is defined as the production of a
re-runnable artefact against an open interface, not a self-graded checklist;
the artefact is what a third party verifies.
\end{answerbox}

% ----------------------------------------------------------------------------
\section{Observability and audit load}
\label{sec:arg-observability}

\begin{questionbox}[O7.1 --- ``Mandating observability everywhere creates a surveillance and cost burden.'']
The observability layer asks for logs, traces, and metrics across the estate.
That is a privacy hazard, a storage cost explosion, and an alert-fatigue
generator. The audit load you create may be worse than the threats you
mitigate.
\end{questionbox}

\begin{answerbox}[Response O7.1]
\textbf{Concession.} Conceded that \emph{indiscriminate collection is harmful}.
Collecting everything is a privacy hazard, a storage liability, and a path to
alert fatigue --- and over-collected personal data is itself a compliance and
breach-exposure risk, not a control.

\medskip
\textbf{Rebuttal.} The objection conflates \emph{indiscriminate} collection
with \emph{purpose-bound, classified} observability. The architecture
specifies observability tied to data classification and retention policy:
what is collected, why, for how long, and who may read it are constrained by
design, not maximised (\trref{Observability and audit}). Auditability is a
legal obligation that exists independently of this document; the architecture
discharges it \emph{proportionately} rather than maximally.

\begin{itemize}[leftmargin=1.4em,itemsep=2pt]
  \item \evlegal Security of processing and breach-detection duties under
        \institution{GDPR} \citep{gdpr} and the technical risk-management
        requirements of \institution{NIS2} \citep{nis2} \emph{require}
        adequate logging --- with the Commission's technical
        articulation for digital-sector entities applied here by
        analogy \citep{eu-regulation-2024-2690};
        the same legal frame caps it through data minimisation and purpose
        limitation \citep{gdpr}. Audit is mandated; over-collection is
        prohibited --- proportionality is the lawful design point.
  \item \evstd Detect-and-respond functions are baseline controls
        \citep{nist-csf,nist-sp800-53,iso27001}; structured, standard
        telemetry (\institution{syslog}, \institution{IPFIX},
        \institution{OpenTelemetry}) makes collection efficient and bounded
        rather than sprawling
        \citep{rfc-5424-syslog,rfc-7011-ipfix,opentelemetry-spec}.
  \item \evempirical Breach economics reward exactly this: faster detection
        and containment measurably reduce breach cost --- breaches
        contained faster cost measurably less, the breach lifecycle
        being a major cost driver
        \citep{ibm-cost-breach-2024,verizon-dbir2024}. Proportionate
        observability is the highest-return control, not overhead.
  \item \evinst Research universities run classification-driven observability
        in production through their security operations and incident-response
        capabilities \citep{umich-safecomputing,stanford-uit}, and the
        sector operates shared trust and incident infrastructure
        (e.g.\ \institution{REN-ISAC} in the US, \institution{FIRST} globally) on this basis
        \citep{ren-isac,first}.
\end{itemize}

\textbf{Bottom line.} The mandate is proportionate, classified observability
--- legally required to exist and legally bounded in extent --- not blanket
surveillance; the standard telemetry interfaces keep its cost bounded.
\end{answerbox}

% ----------------------------------------------------------------------------
\section{Reversibility}
\label{sec:arg-reversibility}

\begin{questionbox}[O8.1 --- ``Reversibility is a comforting promise that never survives contact with production.'']
You promise that each wave is reversible. In practice, once identity,
network, and policy changes are live, rolling back is a fantasy --- the
dependencies entangle, the data has moved, and nobody dares revert. The
reversibility playbook is reassurance, not engineering.
\end{questionbox}

\begin{answerbox}[Response O8.1]
\textbf{Concession.} Conceded that \emph{reversibility is not free and is
often claimed without being engineered}. A rollback that was never designed,
rehearsed, or tested will indeed fail under pressure, and entangled
dependencies can make a naive revert impossible.

\medskip
\textbf{Rebuttal.} That is precisely why reversibility is treated as a
\keyword{designed and rehearsed property}, not an aspiration. Each wave
defines its rollback preconditions, its data-preservation guarantees, and a
tested revert path \emph{before} the wave is committed; coexistence with the
incumbent during a wave is what keeps the revert path open
(\trref{Reversibility playbook}; \trref{Coexistence patterns library}).
Reversibility here is a fault-tolerance requirement, engineered like any
other.

\begin{itemize}[leftmargin=1.4em,itemsep=2pt]
  \item \evstd Designed recoverability and graceful rollback are core
        cyber-resilience and dependable-systems requirements, not optional
        extras \citep{nist-sp800-160v2,avizienis2004}; resilient systems are
        specified to fail back to a known-good state.
  \item \evlegal Reversibility is also a continuity and exit obligation:
        processor-change and data-portability provisions presuppose that a
        configuration can be unwound and data recovered in usable form
        \citep{gdpr,eu-sccs2021}.
  \item \evsector Change managed with explicit rollback gates is standard IT
        governance practice \citep{cobit,dehaes2015}.
  \item \evinst Coexistence-based migrations preserve the revert path in
        production: \institution{Science DMZ} runs alongside the general
        network so the institution can fall back, and federated identity
        overlays --- rather than replaces --- incumbent systems during
        transition
        \citep{science-dmz-2013,esnet-cases,mit-kerberos,umich-its}.
\end{itemize}

\textbf{Bottom line.} Reversibility is engineered and rehearsed per wave, and
kept open by coexistence with the incumbent --- it is a fault-tolerance
requirement, not a comforting promise.
\end{answerbox}

% ============================================================================
\part{Quick-Reference Summary}
% ============================================================================

\section{Objection map}
\label{sec:arg-summary}

The table below summarises every objection in this module, its concession,
the dominant evidence layers, and the one-line rebuttal. \evlegal{}Legal,
\evstd{}Standards, \evsector{}Sector, \evempirical{}Empirical, and
\evinst{}Institutional mark the converging evidence layers used in each
response (minimum three per response).

\renewcommand{\arraystretch}{1.35}
\begin{longtable}{@{}p{0.06\textwidth}p{0.20\textwidth}p{0.20\textwidth}p{0.13\textwidth}p{0.29\textwidth}@{}}
\toprule
\textbf{ID} & \textbf{Objection} & \textbf{Conceded core} &
\textbf{Evidence} & \textbf{One-line rebuttal} \\
\midrule
\endfirsthead
\toprule
\textbf{ID} & \textbf{Objection} & \textbf{Conceded core} &
\textbf{Evidence} & \textbf{One-line rebuttal} \\
\midrule
\endhead
\midrule
\multicolumn{5}{r}{\footnotesize\itshape continued on next page}\\
\endfoot
\bottomrule
\endlastfoot

O1.1 & Architecture is over-engineered for a university &
Unmatched complexity is a real risk &
Std, Sector, Inst &
Layers name functions a campus already runs; contracts \emph{bound} complexity, not add it. \\

O1.2 & Identity/policy are single points of failure; an SLA beats this &
Mediation-point risk is real and architecture-created &
Std, Inst &
Redundancy per plane, PEP caching over token lifetime, declared per-zone fail modes --- engineered, measured, rehearsed. \\

O2.1 & Open-source MDM cannot manage iOS/Android &
No mature open mobile MDM at scale (true) &
Legal, Std, Sector, Emp, Inst &
Open-foundations means contract-bounded replaceability; a conforming proprietary MDM is a valid instantiation. \\

O2.2 & Open source is itself the supply-chain risk &
Abandonware is real; a licence confers no fix &
Std, Sector, Inst &
Support subscriptions where SLAs matter, SLSA provenance, Operational-dimension scoring, replaceability as backstop. \\

O3.1 & Sizing is fantasy for small institutions &
One-size instantiation would be unserious &
Legal, Std, Sector, Inst &
Minimal functional set is size-invariant (no-threshold law); instantiation scales to shared/managed services. \\

O4.1 & Greenfield design; no rip-and-replace is possible &
Rip-and-replace would fail &
Std, Sector, Emp, Inst &
Playbook is brownfield by construction: waves coexist with incumbents and are independently valuable. \\

O5.1 & Open stack just swaps one lock-in for another &
Open licensing alone does not defeat lock-in &
Legal, Std, Sector, Inst &
Replaceability comes from four conditions + conformance tests, not from the logo. \\

O6.1 & Conformance tests are a checklist, not a guarantee &
An unenforced checklist is worthless &
Std, Sector, Emp, Inst &
Conformance = a re-runnable artefact against an open interface, third-party verifiable. \\

O7.1 & Observability is surveillance and cost burden &
Indiscriminate collection is harmful &
Legal, Std, Emp, Inst &
Mandate is proportionate, classified observability --- legally required and legally bounded. \\

O8.1 & Reversibility never survives production &
Un-rehearsed rollback fails &
Legal, Std, Sector, Inst &
Reversibility is designed and rehearsed per wave, kept open by coexistence. \\

\end{longtable}

\bigskip
\noindent\textbf{Methodological note.} Every response in this module concedes
the legitimate core of its objection before rebutting the overreach, and
each rests on at least three independent evidence layers. The recurring
defensive structure --- \emph{size-invariant minimal functional set},
\emph{contract-bounded replaceability}, and \emph{brownfield coexistence} ---
is what allows the architecture to be simultaneously rigorous and applicable
across institutions of very different scale and jurisdiction.

% ============================================================================
\backmatter
% ============================================================================

\printbibliography[heading=bibintoc]

\end{document}
